% GENERATED FILE. Edit canonical lessons, then run npm run generate:books.
% canonical-source-hash: fnv1a64:7c3cbad44e542d8e
% canonical-lessons: KA-C266-balb, KA-C266-rimot, KA-C266-madake, KA-C266-karpura, KA-C266-tamra, KA-R266-first-pass-words-from-chapters-262-263, KA-R266-second-pass-words-from-chapters-264-266

\chapter{Light bulb, Remote control, Earthen pot, Camphor, Copper}
\label{ch:ka-light-bulb-remote-control-earthen-pot-camphor-copper}
\hlchaptermodality{266}

\begin{chapteropening}
\textbf{By the end of this chapter:} \emph{I can say a light bulb, a remote control, an earthen pot, camphor and copper.}

Complete the last lesson of chapter 266: I can say a light bulb, a remote control, an earthen pot, camphor and copper.
\end{chapteropening}

\section[\texorpdfstring{\kn{ಬಲ್ಬ್} (balb)}{ಬಲ್ಬ್ (balb)}]{\kn{ಬಲ್ಬ್} (balb) — a light bulb}

\label{lesson:KA-C266-balb}



\begin{quote}
Before the new one: say the Kannada for a packet, then the Kannada for steel.
\end{quote}

\subsection*{You'll want to know: \kn{ಬಲ್ಬ್}}
\textbf{\kn{ಬಲ್ಬ್}} — \emph{balb} — \textquotedblleft{}a light bulb\textquotedblright{}.

\begin{grammarlens}[title={What you've built}]
One more word for this chapter.
\end{grammarlens}

\subsection*{Your turn}
\begin{itemize}
\raggedright
  \item \emph{Say it:} \emph{balb}
  \item \emph{Say it:} \emph{balb}, once more
  \item \emph{Say it:} say \emph{balb}
  \item \emph{Recall:} say the Kannada for a packet, then the Kannada for steel, then say \emph{balb} again
\end{itemize}

\begin{tcolorbox}[breakable,colback=teal!4,colframe=teal!35!black,title={Before you move on}]
What does \kn{ಬಲ್ಬ್} mean? (\textquotedblleft{}A light bulb\textquotedblright{}.) Say it once more.
\end{tcolorbox}
\section[\texorpdfstring{\kn{ರಿಮೋಟ್} (rimōṭ)}{ರಿಮೋಟ್ (rimōṭ)}]{\kn{ರಿಮೋಟ್} (rimōṭ) — a remote control}

\label{lesson:KA-C266-rimot}



\begin{quote}
Before the new one: say the Kannada for steel, then the Kannada for a light bulb.
\end{quote}

\subsection*{You'll want to know: \kn{ರಿಮೋಟ್}}
\textbf{\kn{ರಿಮೋಟ್}} — \emph{rimōṭ} — \textquotedblleft{}a remote control\textquotedblright{}.

\begin{grammarlens}[title={What you've built}]
One more word for this chapter.
\end{grammarlens}

\subsection*{Your turn}
\begin{itemize}
\raggedright
  \item \emph{Say it:} \emph{rimōṭ}
  \item \emph{Say it:} \emph{rimōṭ}, once more
  \item \emph{Say it:} say \emph{rimōṭ}
  \item \emph{Recall:} say the Kannada for steel, then the Kannada for a light bulb, then say \emph{rimōṭ} again
\end{itemize}

\begin{tcolorbox}[breakable,colback=teal!4,colframe=teal!35!black,title={Before you move on}]
What does \kn{ರಿಮೋಟ್} mean? (\textquotedblleft{}A remote control\textquotedblright{}.) Say it once more.
\end{tcolorbox}
\section[\texorpdfstring{\kn{ಮಡಕೆ} (maḍake)}{ಮಡಕೆ (maḍake)}]{\kn{ಮಡಕೆ} (maḍake) — an earthen pot}

\label{lesson:KA-C266-madake}



\begin{quote}
Before the new one: say the Kannada for a light bulb, then the Kannada for a remote control.
\end{quote}

\subsection*{You'll want to know: \kn{ಮಡಕೆ}}
\textbf{\kn{ಮಡಕೆ}} — \emph{maḍake} — \textquotedblleft{}an earthen pot\textquotedblright{}.

\begin{grammarlens}[title={What you've built}]
One more word for this chapter.
\end{grammarlens}

\subsection*{Your turn}
\begin{itemize}
\raggedright
  \item \emph{Say it:} \emph{maḍake}
  \item \emph{Say it:} \emph{maḍake}, once more
  \item \emph{Say it:} say \emph{maḍake}
  \item \emph{Recall:} say the Kannada for a light bulb, then the Kannada for a remote control, then say \emph{maḍake} again
\end{itemize}

\begin{tcolorbox}[breakable,colback=teal!4,colframe=teal!35!black,title={Before you move on}]
What does \kn{ಮಡಕೆ} mean? (\textquotedblleft{}An earthen pot\textquotedblright{}.) Say it once more.
\end{tcolorbox}
\section[\texorpdfstring{\kn{ಕರ್ಪೂರ} (karpūra)}{ಕರ್ಪೂರ (karpūra)}]{\kn{ಕರ್ಪೂರ} (karpūra) — camphor}

\label{lesson:KA-C266-karpura}



\begin{quote}
Before the new one: say the Kannada for a remote control, then the Kannada for an earthen pot.
\end{quote}

\subsection*{You'll want to know: \kn{ಕರ್ಪೂರ}}
\textbf{\kn{ಕರ್ಪೂರ}} — \emph{karpūra} — \textquotedblleft{}camphor\textquotedblright{}.

\begin{grammarlens}[title={What you've built}]
One more word for this chapter.
\end{grammarlens}

\subsection*{Your turn}
\begin{itemize}
\raggedright
  \item \emph{Say it:} \emph{karpūra}
  \item \emph{Say it:} \emph{karpūra}, once more
  \item \emph{Say it:} say \emph{karpūra}
  \item \emph{Recall:} say the Kannada for a remote control, then the Kannada for an earthen pot, then say \emph{karpūra} again
\end{itemize}

\begin{tcolorbox}[breakable,colback=teal!4,colframe=teal!35!black,title={Before you move on}]
What does \kn{ಕರ್ಪೂರ} mean? (\textquotedblleft{}Camphor\textquotedblright{}.) Say it once more.
\end{tcolorbox}
\section[\texorpdfstring{\kn{ತಾಮ್ರ} (tāmra)}{ತಾಮ್ರ (tāmra)}]{\kn{ತಾಮ್ರ} (tāmra) — copper}

\label{lesson:KA-C266-tamra}



\begin{quote}
Before the new one: say the Kannada for an earthen pot, then the Kannada for camphor.
\end{quote}

\subsection*{You'll want to know: \kn{ತಾಮ್ರ}}
\textbf{\kn{ತಾಮ್ರ}} — \emph{tāmra} — \textquotedblleft{}copper\textquotedblright{}.

\begin{grammarlens}[title={What you've built}]
One more word for this chapter.
\end{grammarlens}

\subsection*{Your turn}
\begin{itemize}
\raggedright
  \item \emph{Say it:} \emph{tāmra}
  \item \emph{Say it:} \emph{tāmra}, once more
  \item \emph{Say it:} say \emph{tāmra}
  \item \emph{Recall:} say the Kannada for an earthen pot, then the Kannada for camphor, then say \emph{tāmra} again
\end{itemize}

\begin{tcolorbox}[breakable,colback=teal!4,colframe=teal!35!black,title={Before you move on}]
What does \kn{ತಾಮ್ರ} mean? (\textquotedblleft{}Copper\textquotedblright{}.) Say it once more.
\end{tcolorbox}
\section[(dialogue)]{Words from chapters 262-263 — a first pass}

\label{lesson:KA-R266-first-pass-words-from-chapters-262-263}



\begin{quote}
Say the words for camphor and copper.
\end{quote}

\subsection*{Your turn}
Say these aloud:

\begin{itemize}
\raggedright
  \item a festoon of mango leaves, kumkum, holige, food offered to a god and then shared, jewellery
  \item a boy, a girl, a man, a woman, an old man
\end{itemize}

\begin{tcolorbox}[breakable,colback=teal!4,colframe=teal!35!black,title={Before you move on}]
A festoon of mango leaves? (\textbf{\kn{ತೋರಣ}}, \emph{tōraṇa}.) Kumkum? (\textbf{\kn{ಕುಂಕುಮ}}, \emph{kuṅkuma}.) Holige? (\textbf{\kn{ಹೋಳಿಗೆ}}, \emph{hōḷige}.) Food offered to a god and then shared? (\textbf{\kn{ಪ್ರಸಾದ}}, \emph{prasāda}.) Jewellery? (\textbf{\kn{ಆಭರಣ}}, \emph{ābharaṇa}.) A boy? (\textbf{\kn{ಹುಡುಗ}}, \emph{huḍuga}.) A girl? (\textbf{\kn{ಹುಡುಗಿ}}, \emph{huḍugi}.) A man? (\textbf{\kn{ಗಂಡಸು}}, \emph{gaṇḍasu}.) A woman? (\textbf{\kn{ಹೆಂಗಸು}}, \emph{heṅgasu}.) An old man? (\textbf{\kn{ಮುದುಕ}}, \emph{muduka}.)
\end{tcolorbox}
\section[(dialogue)]{Words from chapters 264-266 — a second pass}

\label{lesson:KA-R266-second-pass-words-from-chapters-264-266}



\begin{quote}
Say the words for the internet and a website.
\end{quote}

\subsection*{Your turn}
Say these aloud:

\begin{itemize}
\raggedright
  \item the internet, a website, a password, an app, a video
  \item a tin, a water pitcher, a jar, a packet, steel
  \item a light bulb, a remote control, an earthen pot, camphor, copper
\end{itemize}

\begin{tcolorbox}[breakable,colback=teal!4,colframe=teal!35!black,title={Before you move on}]
The internet? (\textbf{\kn{ಇಂಟರ್ನೆಟ್}}, \emph{iṇṭarneṭ}.) A website? (\textbf{\kn{ವೆಬ್ಸೈಟ್}}, \emph{vebsaiṭ}.) A password? (\textbf{\kn{ಪಾಸ್ವರ್ಡ್}}, \emph{pāsvarḍ}.) An app? (\textbf{\kn{ಆಪ್}}, \emph{āp}.) A video? (\textbf{\kn{ವಿಡಿಯೋ}}, \emph{viḍiyō}.) A tin? (\textbf{\kn{ಡಬ್ಬಿ}}, \emph{ḍabbi}.) A water pitcher? (\textbf{\kn{ಬಿಂದಿಗೆ}}, \emph{bindige}.) A jar? (\textbf{\kn{ಜಾಡಿ}}, \emph{jāḍi}.) A packet? (\textbf{\kn{ಪೊಟ್ಟಣ}}, \emph{poṭṭaṇa}.) Steel? (\textbf{\kn{ಉಕ್ಕು}}, \emph{ukku}.) A light bulb? (\textbf{\kn{ಬಲ್ಬ್}}, \emph{balb}.) A remote control? (\textbf{\kn{ರಿಮೋಟ್}}, \emph{rimōṭ}.) An earthen pot? (\textbf{\kn{ಮಡಕೆ}}, \emph{maḍake}.) Camphor? (\textbf{\kn{ಕರ್ಪೂರ}}, \emph{karpūra}.) Copper? (\textbf{\kn{ತಾಮ್ರ}}, \emph{tāmra}.)
\end{tcolorbox}
